\chapter{Conclusion}
\label{ch:conclusion}

This thesis addressed one limitation of the software project management platforms in wide use today: because a single provider stores and controls the history of a project, nobody else can verify that the history has not been tampered with. Starting from the architectural pattern validated for decentralized data management in the construction industry~\cite{adel2023}, the thesis adapted that pattern to software project management and implemented it in a working prototype, ChainBoard.

ChainBoard models accounts, projects and tasks as chaincode-governed state on a permissioned Hyperledger Fabric network. Its operations for creating projects, managing membership and roles, assigning tasks, changing status, commenting and archiving are all recorded in the immutable transaction history of Fabric, not in a log controlled by a provider. The chaincode itself decides who may do what, from the roles stored in the project, so the permission rules do not depend on the backend or the user interface. File attachments are stored on IPFS and anchored to the ledger by their CIDs, which keeps binary data off the chain while still tying each recorded reference to the exact file content. A backend API and a browser frontend offer these functions through a task board similar to Trello or Jira. The audit trail, the main guarantee of the system, is visible and can be checked: every change shows who made it and when, and a verification step checks each record against the ledger and produces a downloadable report with a digest chain that can be verified again offline.

The experiments of \Cref{ch:results} support these claims with concrete results. All eleven functional scenarios passed on a real Fabric network; every forbidden action was refused by the chaincode, also when the interface was bypassed; every manipulation of an audit report was detected; and all 278 automated tests pass. The cost of the guarantee is about two seconds per write, while a read takes about 30 milliseconds. In the user study, all ten participants completed every task and the interface scored 86 on the System Usability Scale. Overall, the Fabric-plus-IPFS pattern carries over to software project management with the same architecture; what changes is the data model and the workflow offered to users.

\section{Limitations}
\label{sec:limitations}

The prototype was evaluated on a local Fabric network with two organizations, not on a geographically distributed or production-scale deployment. The performance figures describe one machine under a synthetic load, not realistic use by many people, and they include no comparison with a conventional database. The most important limitation is the identity model. The Fabric gateway authenticates the organization and not each person, so the backend authenticates people with its own login sessions and tells the chaincode who is acting. The chaincode enforces the \emph{rules} but has to trust the backend about \emph{who} is acting, and a compromised backend could claim any identity. Reads are not restricted by role, so any signed-in person who knows a project identifier can read it. Data written to the ledger is permanent, which conflicts with the right to erasure of the GDPR~\cite{gdpr} for personal data such as names and comments. Attachments are anchored by CID, but nothing in the prototype guarantees that the files stay available. Finally, the data model is deliberately small. It supports a Kanban-style board but not the Scrum framework (there is no backlog ordering, no sprints, no estimation and no burndown chart), and it does not yet cover sub-tasks, dependencies between tasks, notifications or access control per field.

\section{Future Work}
\label{sec:future-work}

The limitations above suggest several directions. The most useful is to give every person their own Fabric identity, with a certificate issued by the Fabric certificate authority, so that the chaincode reads the caller from the client identity of the transaction instead of trusting the backend, as the Fabric-based tool of \textcite{batista2021} already does, and to require endorsement by both organizations for critical actions such as archiving a project. Attribute-based policies of the kind studied by \textcite{gordijn2022combining} could then express the project roles directly in the certificates. Restricting reads by role, or placing sensitive fields in Fabric private data collections and keeping personal data off-chain with only hashes on the ledger, would address confidentiality and erasure. Evaluating the system on a geographically distributed network with a realistic number of organizations and peers, and comparing it with a conventional database, would show its throughput and latency under conditions closer to production. The user study showed that participants value a verifiable history but do not always understand the result of the manual check, so the check could run automatically in the background and show only its outcome, without technical terms; a longer study with real teams over several weeks would also show how the tool is used in practice. Extending the data model with sprints, a backlog, story points, burndown charts, sub-tasks, dependencies and notifications would close the functional gap with mature centralized platforms while keeping the same audit guarantees, and pinning attachments, with a check of the CID on download, would make file storage as dependable as the ledger. Finally, the architectural pattern is not specific to software projects, and applying it to other collaborative domains that need a verifiable history would show how far it generalizes.
